This paper presents a corpus of 141,013 public Slovene-language posts by 432 authors on Bluesky, spanning August 2023 to mid-April 2026. The corpus was built using a protocol-native ATProto workflow combining relay-based discovery of Slovene‑tagged posts, DID‑resolved backfill of surviving public post histories from authors’ Personal Data Servers, and post‑level language filtering validated through manual annotation.
As one of the first language-specific empirical studies of Bluesky use for a small European language community, the corpus documents the emergence of a new public space on an open‑protocol platform. Based on the collected data, a marked increase in posting activity is visible in late 2024. Across the collection period, reply posts account for a large share of the material (68.3\%), including in the months before this increase, indicating that conversational use was already present prior to the surge. Further observed characteristics include sparse and largely account‑concentrated hashtag use, and content sharing shaped primarily by Slovene news media, GIF‑based informal exchange, and Bluesky‑native tooling.
The paper illustrates how protocol‑level access in ATProto enables the construction of a high‑precision corpus for an under‑resourced language without relying on platform‑specific research APIs. At the same time, it raises legal and ethical questions that arise when small and highly identifiable language communities become reconstructable from public, decentralised infrastructures. The paper documents the collection workflow, its validation, and the basic temporal, interactional, and linking properties of Slovene‑language activity on Bluesky, offering a reproducible case study for digital humanities research on open‑protocol social media.